\documentclass[12pt]{article}
\usepackage{hyperref}

\usepackage{graphicx}

\usepackage{mathtools}
\DeclarePairedDelimiter\ceil{\lceil}{\rceil}
\DeclarePairedDelimiter\floor{\lfloor}{\rfloor}

\usepackage{cancel}

\usepackage{amsmath}
\usepackage{amsfonts}
\usepackage{amssymb}

\usepackage{xcolor}
\usepackage{listings}
\definecolor{darkgreen}{rgb}{0.0, 0.2, 0.13}
\lstset{language=C++,
                basicstyle=\ttfamily,
                keywordstyle=\color{blue}\ttfamily,
                stringstyle=\color{red}\ttfamily,
                commentstyle=\color{darkgreen}\ttfamily,
                morecomment=[l][\color{magenta}]{\#}
}

\usepackage{breqn}
\usepackage[]{algorithm2e}
\usepackage{physics}

\newcommand\fnurl[2]{%
  \href{#2}{#1}\footnote{\url{#2}}%
}

\newcommand{\Four}{\mathcal{F}}
\renewcommand{\(}{\left(}
\renewcommand{\)}{\right)}
%% \renewcommand{\{}{\left{}
%% \renewcommand{\}}{\right}}

\usepackage{natbib}
\usepackage{bibentry}
\bibliographystyle{plain}

\title{Multi-GPU design plan for rocFFT}
\author{}



\def\no#1{\nobreak{#1}} % stop equation breaking in dmath

\begin{document}
\maketitle

\section{Introduction}

These are design notes for multi-device FFTs with rocFFT.

\section{Notable use cases}

\subsection{Molecular simulators}

\subsubsection{lammps}
In
\url{https://github.com/lammps/lammps/blob/develop/src/KSPACE/pppm.cpp}
\begin{lstlisting}
void PPPM::compute(...) {
   // ...
   brick2fft();
   poisson();
   // ...
}
\end{lstlisting}
The function \lstinline{brick2fft} is defined in the same file; it
takes bricks of arbitrary local dimension and then called
\lstinline{Remap::perform} which calles \lstinline{remap_3d} in
\href{https://github.com/lammps/lammps/blob/develop/src/KSPACE/remap.cpp}{remap.cpp}
which creates a pencil (ie 2D) decomposition of the data over over the
given MPI procs.

\subsubsection{gromacs}
Gromacs has a slightly different behaviour, where is seems that when
using, for example, 4 GPUs, then data starts distributed between the
GPUs, but is gathered onto a single GPU for spectral operations, and
then presumablky this is presumably scattered to the GPUs.  This is
either done in an MPI setting, or a process-local setting, depending
on the runtime mode for gromacs.

\subsection{Large-scale fluid solvers}

GESTS, CHOLLA

These are partial differential equation solvers.  GESTS looks at high
Reynolds-number turbulence, and cholla looks at interstellar fluid
simulations.  These both make use of very large FFTs.  For GESTS, this
is definitely the critical performance stage.

\section{Existing implementations}

\begin{itemize}
  \item \texttt{vkFFT} doesn't yet support multi-gpu.
  \item \texttt{MKL}: not sure about multi-device support.
\end{itemize}
    
\subsection{Single-node multi-GPU}
\subsubsection{cuFFTXt}
From the v12
\href{https://docs.nvidia.com/cuda/cufft/index.html#multiple-gpu-cufft-transforms}{documentation}
and
\href{https://docs.nvidia.com/cuda/cufft/index.html#cufft-code-examples}{examples}:
\begin{itemize}
\item Up to 16 GPUs
\item 1D:
  \begin{itemize}
  \item only supports power-of-two sizes and power-of-two number of
    GPUs
  \item only complex-to-complex
  \end{itemize}
\item 2D and 3D
  \begin{itemize}
  \item real/complex: requires even-length for r/c, only supports
    in-place, shuffled output.
  \item single-transform has permuted output
  \item you can also just call multiple independent transforms
  \item no Bluestein
  \item single-transform multi-GPU output is permuted.
  \end{itemize}
\end{itemize}
The devices are associted with a plan via
\begin{lstlisting}
cufftXtSetGPUs(cufftHandle plan,
               int nGPUs,
               int *whichGPUs);
\end{lstlisting}
Execution is done by \lstinline{cufftXtExecDescriptorC2C()} (or the
real/complex, single/double version) as described
\href{https://docs.nvidia.com/cuda/cufft/#multiple-gpu-2d-and-3d-transforms-on-permuted-input}{here}.
These take \lstinline{cudaLibXtDesc *input} and output arguments.
These structs are allocated with
\begin{lstlisting}
cufftXtMalloc(cufftHandle plan, 
              cudaLibXtDesc ** descriptor,
              cufftXtSubFormat format);
\end{lstlisting}
Of note is that the application doesn't seem to have direct access to
the gpu buffer; the \lstinline{cudaLibXtDesc} is created with
\lstinline{cufftXtMalloc}, and then we use these descriptor structs in
the call to execute the transform.  This seems like a huge
disadvantage - do applications need to copy the data to the new
buffer?



\subsection{Multi-node}

A
\href{https://icl.utk.edu/files/publications/2022/icl-utk-1548-2022.pdf}{technical
  report} is available from ICL and the University of Tennessee.  This
compares a AccFFT, 2Decomp\&FFT, FFTE, FFTW, FFTMPI, heFFTe, SWFFT,
P3DFFT, and FFTADVMPI.  It provides the results of benchmark tests
on summit and spock.

\subsubsection{heFFTe}

\href{heffte}{https://github.com/icl-utk-edu/heffte} is an open-source
DOE-sponsored project.  It provides pencil and slab decomposition, AMD
and NVIDIA back-ends.


\subsubsection{AccFFT}

\href{http://accfft.org/}{AccFFT} supports cuFFT as a back-end.
Pencil and slab decomposition are supported, and spectral operators
are mentioned.

\subsubsection{FFTE}

\url{https://github.com/certik/ffte} is a Fortran code on CPU/Nvidia
GPUs.

\subsubsection{CPU codes}

FFTW has an MPI implimentation, which provides a 1D data
decomposition, ie slab.  Interesting algorithm: round-robin based
send/receive.  Online tuning framework determines the best transpose
algorithm.  The local data distribution is determined by
FFTW\footnote{\url{http://fftw.org/fftw3_doc/MPI-Data-Distribution.html}}:
\begin{lstlisting}
alloc_local
= fftw_mpi_local_size_2d(N0, N1,
                        MPI_COMM_WORLD,
                        &local_n0, &local_0_start);
\end{lstlisting}

p3dFFT \url{https://github.com/sdsc/p3dfft} handles pencil.
\url{https://github.com/sdsc/p3dfft.3}{P3DFFT++}.  Mentions a 3D
domain decomposition (coming soon)?  The interface is provided in
\href{https://github.com/sdsc/p3dfft/blob/master/include/p3dfft.h}{p3dfft.h}.
There is a GPU-FFT branch at
\href{https://github.com/sdsc/p3dfft.3/tree/cuda-dev}{cuda-dev}.

FFTW++: \url{https://fftwpp.sourceforge.net/},
\url{https://github.com/dealias/fftwpp}.  Pencil and slab
decomposition.  Recursive transpose, which helps with latency-bound
tranposes
(\url{https://malcolmroberts.github.io/publications/transpose0.pdf}{conference
  proceedings}).  Interface is at
\href{https://github.com/dealias/fftwpp/tree/master/mpi/mpifftw++.h}{\texttt{mpifftw++.h}}


\subsubsection{FFTX}

\href{http://www.spiral.net/software/fftx.html}{Spiral/FFTX} generate
kernels for FFTs including spectral operations.  There is some support
for MPI, but details are unclear at this point.

\subsubsection{cuFFTMp}

\href{https://docs.nvidia.com/hpc-sdk/cufftmp/index.html}{cuFFTMp} has
recently been made generally accessible by NVIDIA.  A recent
\href{https://developer.nvidia.com/blog/massively-improved-multi-node-nvidia-gpu-scalability-with-gromacs/}{blog
  post from NVIDIA} shows a good speedup for gromacs, claiming a
2$\times{}$ speedup over heFFTe.  Restrictions are similar to cuFFTXt;
sometimes more severe, sometimes less so.  Multiprocess GPUs with more
than one GPU per process doesn't seem possible.  cuFFTMp is based on
nvshmem, but can interoperate with MPI.

Initialization is either with nvshmem or an MPI comm via the
\href{https://docs.nvidia.com/hpc-sdk/cufftmp/api/index.html#c.cufftMpAttachComm}{function}
\begin{lstlisting}
cufftMpAttachComm(cufftHandle plan,
                  cufftMpCommType comm_type,
                  void *comm_handle);
\end{lstlisting}

One can choose a preset data format via
\href{https://docs.nvidia.com/hpc-sdk/cufftmp/api/index.html#cufftxtsetsubformatdefault}{cufftxtsetsubformatdefault} which has format
\begin{lstlisting}
cufftXtSetSubformatDefault(cufftHandle plan,
                           cufftXtSubFormat subformat_forward,
                           cufftXtSubFormat subformat_inverse)
\end{lstlisting}
where the \lstinline{cufftXtSubFormat}s are taken from the enum values
\begin{lstlisting}
CUFFT_XT_FORMAT_INPLACE
CUFFT_XT_FORMAT_INPLACE_SHUFFLED
CUFFT_XT_FORMAT_DISTRIBUTED_INPUT
CUFFT_XT_FORMAT_DISTRIBUTED_OUTPUT
\end{lstlisting}
These are slab input/output for the first two and box inputs in the
latter two as described in
\href{https://docs.nvidia.com/hpc-sdk/cufftmp/api/index.html#c.cufftXtSubFormat}{cufftXtSubFormat}.

From
\href{https://docs.nvidia.com/hpc-sdk/cufftmp/usage/api_usage.html#usage-with-custom-slabs-and-pencils-data-decompositions}{the
  documentation}, ``cuFFTMp also supports arbitrary data distributions
in the form of 3D boxes.''  General arbitrary box distributions are
set via
\href{https://docs.nvidia.com/hpc-sdk/cufftmp/api/index.html#cufftxtsetdistribution}{cufftxtsetdistribution}:
\begin{lstlisting}
cufftXtSetDistribution(cufftHandle plan,
                       int rank,
                       const long long int *lower_input,
                       const long long int *upper_input,
                       const long long int *lower_output,
                       const long long int *upper_output,
                       const long long int *strides_input,
                       const long long int *strides_output)
\end{lstlisting}
where \lstinline{int rank} is the rank of the transform (ie not the
MPI rank or the equivalent with nvshmem).  Execution is done via
\href{https://docs.nvidia.com/hpc-sdk/cufftmp/api/index.html#c.cufftXtExec}{cufftXtExec}
which has signature
\begin{lstlisting}
cufftXtExec(cufftHandle plan,
            void *idata,
            void *odata,
            int direction);
\end{lstlisting}
where \lstinline{void *idata} and \lstinline{void *odata} are GPU
global memory pointers.

There is a reshape function with similar functionality, which used
separately from the cuFFTMp library itself.  Differences between the
single-node/multi-gpu and multi-process interface is given
\href{https://docs.nvidia.com/hpc-sdk/cufftmp/usage/differences.html}{here}.

\section{Design objectives}

Gromacs and lammps offer interesting use cases where we could design
and API which would enable us to increase application performance
while making application developer's lives easier.

\subsection{cuFFTXt features/anti-features}
The cuFFTXt interface imposes data distributions on the user
(including restricting problem decomposition) which would seem to
reduce the usability of the library and the performance of the
application code.

Multi-GPU 1D transforms are particularly onerous - the FFT is
bandwidth limited, and this is made worse by going off-device.

Batched multi-GPU transforms are handled independently, with the
distribution of transforms imposed by the cuFFTXt library.  This seems
really onerous (what if the data isn't on the right GPU?) and it's
hard to imagine a use case that is better than the alternative of just
performing a bunch of independent FFTs on different devices.
Moreover, the fact that the data layout changed if the batch size
is/is not equal to one is a bit of a surprise for users.


\subsection{Other ROCm libraries}

\texttt{cuFFTXt} uses \lstinline{cudalibxt.h} which also covers some
\texttt{cuBLAS} functionality.  There are no current plans to add
multi-GPu support to \texttt{rocBLAS}, but it would be good to keep
the possibility in mind so that we can avoid having to back-step on an
API choice.  NVIDIA also made cuSolverMp; if we decide to make
rocSolverMp, it would be good to keep the interfaces as compatible as
possible to avoid pain-points for application developers.

\subsection{Opportunities}

If we are extending the rocFFT API, we have an option to add some
functionality.  In particular, multiple batch dimensions is a feature
that could help.


\section{Proposal ideas}

Question: one buffer per device?

We need to specify, for each device, a starting index. Optional: end
index, stride.  Also for output.

\subsubsection{More specifications}


No struct: 
\begin{lstlisting}
rocfft_status rocfft_plan_description_set_device_layout(
  rocfft_plan_description description,
  size_t device_num,
  size_t in_strides_size,
  size_t* in_strides,
  size_t out_strides_size,
  size_t* out_strides,
  size_t startidx_size,
  size_t* startidx,
  size_t stopidx_size,
  size_t* stopidx);
\end{lstlisting}
which specifies the data layout for device \lstinline{device_num}.
The user loops over the $d$ devices, setting each device's data layout
via separate function call.  The fields are nullable, at which point
the default values are loaded.

We would also need \lstinline{get}er functions for each of these
values.  This is pretty verbose; we would need 

\begin{lstlisting}
rocfft_status rocfft_plan_description_get_device_startidx(
  rocfft_plan plan,
  size_t device_num,
  size_t dimension,
  size_t* startidx}
\end{lstlisting}
for \lstinline{startidx}, \lstinline{endidx}, \lstinline{in_strides},
and \lstinline{out_strides}.

With a struct:

In particular, adding a struct
\begin{lstlisting}
typedef struct {
  size_t istart;
  size_t iend;
  size_t instride;
  size_t outstride;
} rocfft_iodim;
\end{lstlisting}
Which can be passed to the function 
\begin{lstlisting}
rocfft_status rocfft_plan_description_set_device_layout(
  rocfft_plan_description description,
  size_t device_num,
  size_t rocfft_iodim_size,
  rocfft_iodim* rocfft_iodim_s);
\end{lstlisting}

We can access this data by first getting the dimension
\begin{lstlisting}
  size_t rocfft_plan_dimension(... plan);
\end{lstlisting}
and then using that to fill an array of iodims for a given device:
\begin{lstlisting}
rocfft_status rocfft_plan_description_set_device_layout(
  rocfft_plan_description description,
  size_t device_num,
  size_t dimension,
  rocfft_iodim* rocfft_iodims);
\end{lstlisting}

\subsubsection{Computed specifications}

Consider 1D array distributed between $P$ devices.  We are given the
length $\ell$ of the transform.  If there is one device, then we start
at index $0$ and go to $<\ell$.  If we have two devices, then we need
two starting indices, $n_0$ and $n_1$.  Wlog, $n_0=0$.  Suppose that
the stride is 1.  Then, device 0 handles indides $0,\dots,n_1-1$ and
device 1 handles indices $n_1,\dots,\ell-1$.  If the stride is 2, then
$n_0=0,$, $n_1=1$, and device 0 handles indices $0,2,\dots$ and device
1 handles $1, 3, \dots$ (both are bounded by $\ell-1$ - there are
even/odd cases to handle here which we omit for simplicity).

This algorithm can be extended to general indices, per-device strides,
transform strides, and lengths.  Dimensions will act independently, so
this will generalize to arbitrary dimensions.  The data layout can be
checked for aliasing errors by checking per-device data layouts
pairwise using the array correctness algorithm.

Using this method, we can specify the data layout using fewer
parameters.  What's left to specify is
\begin{lstlisting}
typedef struct {
  size_t in_index;
  size_t out_index;
  size_t istride;
  size_t outstride;
} rocfft_iodim;
\end{lstlisting}
for all $d$ dimensions and $P$ devices.  Note that these are
meta-strides; these strides work on the data, which may itself be
strided.  (Question: should we just have index/stride for input and
output?)  This struct is similar to what is used in FFTW's guru
interface, so users may be more familiar with it (and we can also use
this for single-GPU higher-dimension FFTs/batches).  The interface
would be
\begin{lstlisting}
rocfft_status rocfft_plan_description_set_device_layout(
  rocfft_plan_description description,
  size_t device_num,
  size_t rocfft_iodim_size,
  rocfft_iodim* rocfft_iodim_s);
rocfft_status rocfft_plan_get_device_layout(
  rocfft_plan plan,
  size_t device_num,
  size_t dimension,
  rocfft_iodim* rocfft_iodims);
\end{lstlisting}
We should probably think about how to get the ending index from the
plan as well, as this will be very useful for users.  Does this need to
be added to the API?

Note, however, that the order is important here.  Consider the cuFFTMp
example
\href{https://docs.nvidia.com/hpc-sdk/cufftmp/usage/api_usage.html#usage-with-custom-slabs-and-pencils-data-decompositions}{distribution}.
While the Green block would be unambiguous, either the blue or the
purple block would end up with the lower-right quadrant, depending on
which one was processed first.

\subsection{overall behaviour}

We can check whether the parameters are valid at plan generation, and
then return an error code if things aren't set up correctly

\section{API proposal}

\subsection{Single-node multi-gpu}

\texttt{rocFFT} has \lstinline{#if 0}'d-out
\href{https://github.com/ROCmSoftwarePlatform/rocFFT/blob/develop/library/include/rocfft.h#L282}{code}
for specifying the number of devices; this was present in the first
commit, and has never been used.  The API for single-process,
multi-GPU selection is
\begin{lstlisting}
rocfft_plan_description_set_devices(
  rocfft_plan_description description,
  void *devices,
  size_t number_of_devices);
\end{lstlisting}
which seems a perfectly reasonable API. Since the
\lstinline{rocfft_plan_description} is created before passing it to
the plan, we can use this as a vehicle to pass a bunch of information
to the plan generator.

This parameters of the data decomposition can be access via a getter
function.  FIXME: write this

FIXME: how do we set the per-buffer data distribution?

FIXME: how do we exec?  Just pass more in the \texttt{void**}?

FIXME: how do we manage error cases?  TODO: flow chart.

\subsection{Multi-node single-gpu}

FIXME: idea: use the descriptor to attach a communicator if we're
using MPI.

FIXME: idea: use the multi-device layout descriptor to specify that
we're using a 

\subsection{Multi-node multi-gpu}

FIXME: idea: we can attach a communicator as well as a list of local
devices.  Question: how can we figure out if a buffer is associated
with a process?  Looks like nvshmem has its own hyrda launcher; are we
guaranteed something like the MPi rank?

\subsection{hipFFT}

FIXME: how can we get multiple devices per proc via the desciptors?

\bibliography{../refs}


\end{document}
